\section{What two lensed images actually measure}
\label{sec:inference}

This section derives the idealized two-image likelihood used in our sensitivity
calculations. The weak-diffraction signal is a random frequency-dependent
perturbation, as in the framework of \cite{oguri2022,choi2026}; the differential
weights, nuisance projection, and statistical information accounting below are
derived here. The results are sensitivity requirements, not population forecasts.

\subsection{Eliminating a shared waveform and paying its noise cost}

For positive Fourier frequencies, write the two scalar observations as
\begin{equation}
 d_i(f)=\sqrt{|\mu_i|}\,e^{i(2\pi f t_i+\phi_i)}
 h(f)[1+\delta_i(f)]+n_i(f),\qquad i=1,2.
\end{equation}
Here \(h\) is the common unlensed waveform, \(\delta_i\) is a fractional
perturbation, and the independent noises have one-sided power spectral densities
\(S_i(f)\). Dividing out the macro amplitude and phase and differencing removes
the common waveform at zeroth order:
\begin{equation}
 r(f)=h(f)[\delta_1(f)-\delta_2(f)]+n_\Delta(f),\qquad
 S_\Delta(f)=\frac{S_1(f)}{|\mu_1|}+\frac{S_2(f)}{|\mu_2|}.
\end{equation}
Define image SNR densities \(w_i(f)=4|\mu_i||h(f)|^2/S_i(f)\), with
\(\rho_i^2=\int w_i(f)\,df\). The exact differential weight is
\begin{equation}
 \boxed{w_\Delta(f)=\frac{4|h(f)|^2}{S_\Delta(f)}
 =\frac{w_1(f)w_2(f)}{w_1(f)+w_2(f)}.}
 \label{eq:infer_harmonic}
\end{equation}
This also follows by minimizing a two-image Gaussian likelihood over an
arbitrary common complex source amplitude independently in each frequency bin.
Thus source-waveform cancellation has a definite noise cost.

If the two SNR densities have identical normalized frequency dependence,
\(w_i(f)=\rho_i^2p(f)\), where \(\int p(f)\,df=1\), then
\begin{equation}
 \rho_\Delta\equiv\left[\int w_\Delta(f)\,df\right]^{1/2}
 =\frac{\rho_1\rho_2}{\sqrt{\rho_1^2+\rho_2^2}}.
 \label{eq:infer_rhodelta}
\end{equation}
Two equally loud images give \(\rho_\Delta=\rho_i/\sqrt{2}\); for unequal
images the quieter image limits the comparison. Equation~\eqref{eq:infer_rhodelta}
is an upper bound at fixed integrated image SNRs if their spectral shapes differ;
in that case one must first integrate Eq.~\eqref{eq:infer_harmonic}.

This scalar cancellation presumes that detector responses have been accounted
for. At different arrival times a real network measures different mixtures
\(F^+_I(t_i)h_++F^\times_I(t_i)h_\times\). Precession or higher harmonics can
make the polarization ratio frequency dependent. A realistic analysis should
therefore use a joint network likelihood sharing both source polarizations, or
propagate the uncertainty of a polarization reconstruction. A saddle image also
requires the actual Morse/Hilbert-transform operator: its phase is \(-i\) at
positive frequency and \(+i\) at negative frequency, up to Fourier convention.
For general waveforms this is not equivalent to shifting the binary's orbital
coalescence phase. Intrinsic waveform errors can cancel in a consistent joint
comparison; calibration and response errors need their own nuisance model.

\subsection{Projecting out ordinary magnification, phase, and delay}

Macro amplitudes and arrival times are fitted from the same data. Introduce the
real inner product
\begin{equation}
 (u,v)_\Delta=\operatorname{Re}\int df\,w_\Delta(f)u^*(f)v(f)
\end{equation}
and the three nuisance functions \(b_a(f)=\{1,i,2\pi if\}\), with real
coefficients. With \(\Gamma_{ab}=(b_a,b_b)_\Delta\), the retained perturbation is
\begin{equation}
 P_\perp u=u-\sum_{ab}b_a(\Gamma^{-1})_{ab}(b_b,u)_\Delta.
 \label{eq:infer_projection}
\end{equation}
A real rather than complex projection matters: a time shift removes an imaginary
linear slope, not a real amplitude slope. Fitting an arbitrary constant phase is
a conservative choice even if a Morse phase is theoretically discrete.

For \(\Delta\delta=\delta_1-\delta_2\), define the detector-weighted,
projected \emph{differential} fractional RMS by
\begin{equation}
 \epsilon_\Delta^2=\rho_\Delta^{-2}
 \int df\,w_\Delta(f)|P_\perp\Delta\delta(f)|^2.
\end{equation}
For a specified residual realization, its matched-filter SNR is
\begin{equation}
 \boxed{\rho_{\rm res}=\rho_\Delta\epsilon_\Delta.}
 \label{eq:infer_known}
\end{equation}
This is a useful sensitivity ceiling when the realization is not known; it is
not automatically the significance for discovering a stochastic covariance.
For example, two images each with \(\rho_i=500\) give
\(\rho_{\rm res}=3.54\) for \(\epsilon_\Delta=10^{-2}\), or \(0.354\) for
\(10^{-3}\). These are illustrative sensitivity inputs, not predicted CDM
residual amplitudes or event rates.

If each image instead has projected RMS \(\epsilon\), and their weighted real
correlation is \(c\), then
\begin{equation}
 \epsilon_\Delta^2=2\epsilon^2(1-c).
\end{equation}
For equally loud independent images, the extra \(\sqrt{2}\) differential signal
compensates the extra \(\sqrt{2}\) noise, yielding
\(\rho_{\rm res,rms}=\rho_i\epsilon\). A perfectly common perturbation cancels.
The experiment measures differential small-scale structure, and cannot retain
every mode of the total line-of-sight variance.

\subsection{Stochastic detection is a covariance problem}

The complex covariance of the difference is
\begin{equation}
 C_\Delta=C_{11}+C_{22}-C_{12}-C_{21},\qquad
 C_{ij}(f,f')=\langle\delta_i(f)\delta_j^*(f')\rangle.
\end{equation}
One must also retain the nonconjugated covariance
\(J_{ij}(f,f')=\langle\delta_i(f)\delta_j(f')\rangle\), or directly construct
the covariance of the stacked real and imaginary components. A real convergence
field with a complex diffraction kernel does not generally produce a proper,
circular complex process. Our implementation constructs the real covariance,
projects it with Eq.~\eqref{eq:infer_projection}, and whitens it using the
frequency weights. The resulting eigenvalues \(\eta_a\geq0\) are signal-variance
to noise-variance ratios in independent real modes.

On the retained whitened subspace, take \(x\sim N(0,I+\Lambda)\), where
\(\Lambda\) is the lens covariance and \(\eta_a\) are its eigenvalues. Integrating
over a Gaussian random lens realization gives
\begin{equation}
 \ln\frac{p(x|\Lambda)}{p(x|0)}=
 -\frac12\ln\det(I+\Lambda)
 +\frac12x^T[I-(I+\Lambda)^{-1}]x.
\end{equation}
Our covariance-information measure is
\begin{equation}
 \boxed{d\equiv\sqrt{2D_{\rm KL}(p_\Lambda\Vert p_0)}
 =\left\{\sum_a[\eta_a-\ln(1+\eta_a)]\right\}^{1/2}.}
 \label{eq:infer_KL}
\end{equation}
This is \emph{not a calibrated Gaussian discovery significance}. The expected
squared matched-filter SNR if the realization were known is instead
\(\langle\rho_{\rm res}^2\rangle=\sum_a\eta_a\).
For weak covariance, \(\eta_a\ll1\) for all modes,
\begin{equation}
 d^2\simeq \frac12\sum_a\eta_a^2,
\end{equation}
which is also the squared mean shift of the optimal quadratic score in units of
its null standard deviation. At fixed total residual power, spreading that
power over more independent weak modes reduces this covariance discrimination.

Kullback--Leibler divergences add for independent events. Thus \(N\) identical
events have \(d_N=\sqrt{N}\,d_1\). Because weak-signal covariance is proportional
to fractional amplitude squared, \(d\propto\epsilon^2\), the detectable
stochastic RMS amplitude improves only as \(N^{-1/4}\) in this regime. The
\(N^{-1/2}\) amplitude scaling of a known deterministic waveform does not apply
to independent unknown lens realizations. This Gaussian treatment is an
explicit approximation; a sparse, non-Gaussian halo population would require
an appropriate hierarchical likelihood.
